\section{Scope, Method, and Qualifications}

\subsection{Reader, question, and claim boundary}

This is a discursive article in faith and theology for an intellectually serious reader---Christian or considering the Christian claim---who can follow a metaphysical distinction but is not presumed to know scholastic terminology. Its governing question: does the concrete particularity of the biblical economy---one chosen people, one incarnate life, instituted material sacraments, personal address---degrade the universality proper to God, as a long objection tradition from Celsus through Spinoza and Lessing contends? Its governing claim: the objection rests on a false premise that generality is the divine mode; in the classical account, only particulars exist concretely, abstraction is the finite knower's limitation, God knows and sustains every singular directly, and the biblical particulars are consistently ordered to universal scope---election for all nations, one assumed body grounding union with every human being, repeatable concrete signs carrying the once-for-all event outward.

The argument establishes coherence and fittingness, not occurrence. It does not prove the Incarnation or Resurrection from metaphysics, adjudicate among claimed revelations, resolve historical questions such as the Quirinius census date, or supply the reader's act of faith. Fittingness arguments (\emph{convenientia}) are used only retrospectively: they show why an event believed on other grounds is worthy of God, never that it had to happen.

\subsection{Terminology fixed for this article}

\begin{itemize}
\item \textbf{Particular} means concrete, singular, historically located: this person, people, body, act, sign. It is predicated of creatures, divine actions, and the assumed humanity of Christ---never of God as though he were one member of a class.
\item \textbf{Universal} is used in two senses kept explicit in the argument: the abstract universal (a common nature as object of thought) and the universal cause (God as giver of being to every creature). The article's central distinction is that the second does not entail the first as a divine mode of relating to creatures.
\item \textbf{Scandal of particularity} is a received modern label for the objection, not a scriptural or magisterial category.
\item \textbf{Flight into abstraction} names the article's diagnosis of the preference for a non-landing God; it is an editorial characterization, not a canonical classification of any thinker.
\item \textbf{Election} means gratuitous divine choice conferring vocation and mission, expressly not antecedent merit or exclusive divine concern.
\end{itemize}

\subsection{Source hierarchy and authority classes}

Scripture governs the pattern claims (election's gratuity and outward ordering, the Incarnation, the risen Christ's particularity); brief clauses are quoted from the public-domain Douay--Rheims (Challoner) translation as an identified working witness, without any claim that it is a critical text. Conciliar and catechetical texts (\emph{Gaudium et spes} 22; \emph{Catechism of the Catholic Church} 1084--1085) govern the universality of Christ's union with humanity and the sacramental economy, at their own authority level. Patristic witnesses (Irenaeus, Athanasius, Leo) show that the particularity objection was met, not evaded, in the tradition; breadth of citation is illustrative, not an authority count. Aquinas supplies the metaphysical framework (knowledge of singulars, providence over individuals, God outside every genus, the individual term of the assumption, sensible signs); these are school positions adopted and argued, not dogmas. Celsus (as quoted by Origen), Spinoza, and Lessing are objection sources, quoted to state the strongest opposing case, with their governing premises identified rather than conceded. Everything else---the ordering of the argument, the inversion of the abstraction premise, the doctor and form-letter analogies, the reading of Rev 2:17 as particularity perfected---is project synthesis.

\subsection{Material qualifications}

\begin{enumerate}
\item God is not a particular alongside others; predicating choice and address of God does not place him in a genus (ST I, q.~3, a.~5 is cited to that effect).
\item The moderate-realist account of universals and the Thomistic analysis of abstraction are adopted school positions; the theological claims of the later sections do not stand or fall with any one technical formulation of them.
\item Divine knowledge of singulars and particular providence follow the classical Thomistic account; no inference is licensed from particular suffering to particular guilt.
\item Election's outward ordering does not by itself identify which claimed revelation is true; the article says so explicitly and argues coherence only. Israel's election remains, in Christian teaching, irrevocable; nothing here treats Israel as a discarded instrument, and the Israel--Church relation is not adjudicated.
\item Athanasius's cosmic-Word argument is directed \emph{ad hominem} at interlocutors who granted the Word's presence in the whole; the hypostatic union differs in kind from the Word's union with creation, as the cited edition's own annotation notes.
\item ``United Himself in some fashion with every man'' retains the council's qualifier; universal salvific scope is not converted into a claim that all are in fact saved.
\item Sacramental making-present is not repetition of the historical event, and the paragraph on assurance does not reduce grace to mechanism or dispense with faith and disposition; sacramental validity, matter and form, and Eucharistic doctrine are outside scope.
\item The Lessing quotations are given in German from attested transcriptions; the English renderings are the article's own labeled working translations, and no public-domain English translation was identified. No page-identified primary printing was directly inspected; this is recorded as a verification limit in the research record.
\item Celsus survives only in Origen's quotation; attribution is to Celsus as reported.
\item The closing section's claims about love and the white stone are theological reading and synthesis, not exegetical assertions about authorial intention.
\end{enumerate}

\subsection{Coverage, currentness, translations, rights, and review}

No liturgical rite, civil or canonical jurisdiction, or mutable discipline governs any conclusion; no legal claims are made. All online witnesses were checked on 24 July 2026: Douay--Rheims chapter pages, New Advent transcriptions of the Ante-Nicene and Nicene and Post-Nicene Fathers translations and of the English Dominican \emph{Summa}, a full-text NPNF copy of Robertson's Athanasius, the Corpus Thomisticum Latin for ST III, q.~1, a.~1, the Vatican website's English \emph{Catechism} and \emph{Gaudium et spes}, and the Project Gutenberg Elwes translation of Spinoza. Quotations are brief and argument-driven. Scripture, official ecclesiastical texts, and third-party translations retain their own status and are not offered under the project's CC BY 4.0 license. This is an internally source-audited working study by an AI-assisted process; it has received no independent philosophical, biblical, patristic, dogmatic, literary, or ecclesiastical review and claims no ecclesiastical approval.
